\documentclass[../../main.tex]{subfiles}
\begin{document}

{\bf [解答]}

実数$(x,y)$が
\begin{align}
  x^2+xy+y^2=6 \label{eq:1}
\end{align}
を満たす時の
\begin{align}
  f(x,y)
  &= x^2y+xy^2-x^2-2xy-y^2+x+y \\
  &= xy(x+y)-(x+y)^2+(x+y) \label{eq:2}
\end{align}
のとりうる範囲を考える．

$\alpha=x+y$，$\beta=xy$とおくとこれらは$t^2-\alpha t+\beta=0$の二つの実数解であり，判別式$D$として
\begin{align}
  D = \alpha^2-4\beta \ge0 \label{eq:3}
\end{align}
が成り立つ．このもとで\cref{eq:1,eq:2}を$\alpha,\beta$で表すと
\begin{align}
  & \alpha^2-\beta=6 \label{eq:4} \\
  & f(\alpha,\beta)=\beta\alpha-\alpha^2+\alpha \label{eq:5}
\end{align}
となる．

\cref{eq:4}から$\beta=\alpha^2-6$だから\cref{eq:3,eq:5}に代入して
\begin{align}
  &
  \begin{cases}
    \alpha^2-4(\alpha^2-6)\ge0 \\
    f(\alpha)=(\alpha^2-6)\alpha-\alpha^2-5\alpha
  \end{cases} \\
  \therefore
  &
  \begin{cases}
    -2\sqrt{2} \le \alpha \le 2\sqrt{2} \\
    f(\alpha)=\alpha^3-\alpha^2-5\alpha
  \end{cases}
\end{align}
である．$f(\alpha)$の一階微分を計算すると
\begin{align}
  f'(\alpha)=3\alpha^2-2\alpha-5=(3\alpha-5)(\alpha+1)
\end{align}
だから，$f(\alpha)$の増減表は以下のようになる．

\begin{table}[htb]
  \centering
  \caption{$f(\alpha)=\alpha^3-\alpha^2-5\alpha$の$-2\sqrt{2} \le \alpha \le 2\sqrt{2}$での増減表}
  \begin{tabular}{c|ccccccc}
    $\alpha$ & $-2\sqrt{2}$ & & $-1$ & & $5/3$ & & $2\sqrt{2}$ \\
    \hline
    $f'$ & & $+$ & $0$ & $-$ & $0$ & $+$ & \\
    \hline
    $f$ & $-6\sqrt{2}-8$ & $\nearrow$ & $3$ & $\searrow$ & $-\dfrac{175}{27}$ & $\nearrow$ & $6\sqrt{2}-8$ \\
  \end{tabular}
\end{table}

従って，$3>6\sqrt{2}-8$及び$-6\sqrt{2}-8<-\dfrac{175}{27}$に注意して，$f(\alpha)$のとりうる範囲は
\begin{align}
  -2(3\sqrt{2}+4)\le f(\alpha) \le 3
\end{align}
である．

\begin{figure}[htb]
  \centering
  \begin{tikzpicture}[x=1.3cm, y=0.25cm, >=stealth,
    every node/.style={font=\small}]
    \def\twoSqrtTwo{2.828427}
    \def\negTwoSqrtTwo{-2.828427}
    \def\minVal{-16.48528}
    \def\maxVal{3}
    \def\endVal{0.48528}
    \def\localMinVal{-6.48148}

    % 座標軸
    \draw[->, gray!80!black, line width=0.85pt] (-3.5,0) -- (3.6,0) node[right] {$\alpha$};
    \draw[->, gray!80!black, line width=0.85pt] (0,-18.5) -- (0,5.5) node[above] {$y$};

    % 原点
    \node[below left=2pt, font=\footnotesize] at (0,0) {$O$};

    % 補助破線
    % 最小点 (-2\sqrt{2}, -6\sqrt{2}-8)
    \draw[dashed, gray!60, line width=0.7pt]
    (\negTwoSqrtTwo, 0) -- (\negTwoSqrtTwo, \minVal) -- (0, \minVal);
    % 極大点 (-1, 3)
    \draw[dashed, gray!60, line width=0.7pt]
    (-1, 0) -- (-1, 3) -- (0, 3);
    % 極小点 (5/3, -175/27)
    \draw[dashed, gray!60, line width=0.7pt]
    (1.6667, 0) -- (1.6667, \localMinVal) -- (0, \localMinVal);
    % 端点 (2\sqrt{2}, 6\sqrt{2}-8)
    \draw[dashed, gray!60, line width=0.7pt]
    (\twoSqrtTwo, 0) -- (\twoSqrtTwo, \endVal) -- (0, \endVal);

    % 曲線 f(\alpha) = \alpha^3 - \alpha^2 - 5\alpha
    \draw[blue!85!black, line width=1.35pt]
    plot[domain=\negTwoSqrtTwo:\twoSqrtTwo, samples=70] (\x, {\x*\x*\x - \x*\x - 5*\x});

    % 端点・極値の点
    % 最大値 (極大)
    \node[circle, fill=red!85!black, inner sep=1.4pt] at (-1, 3) {};
    \node[above=2pt, red!85!black, font=\footnotesize] at (-1, 3) {最大値 $3$};

    % 最小値 (左端点)
    \node[circle, fill=red!85!black, inner sep=1.4pt] at (\negTwoSqrtTwo, \minVal) {};
    \node[below left=1pt, red!85!black, font=\footnotesize] at (\negTwoSqrtTwo, \minVal) {最小値 $-6\sqrt{2}-8$};

    % 極小点
    \node[circle, fill=blue!85!black, inner sep=1.2pt] at (1.6667, \localMinVal) {};

    % 右端点
    \node[circle, fill=blue!85!black, inner sep=1.2pt] at (\twoSqrtTwo, \endVal) {};

    % 軸目盛りラベル
    \node[above=2pt, font=\footnotesize] at (-1, 0) {$-1$};
    \node[above=2pt, font=\footnotesize] at (\negTwoSqrtTwo, 0) {$-2\sqrt{2}$};
    \node[below=2pt, font=\footnotesize] at (1.6667, 0) {$\frac{5}{3}$};
    \node[below=2pt, font=\footnotesize] at (\twoSqrtTwo, 0) {$2\sqrt{2}$};

    \node[left=2pt, font=\footnotesize] at (0, 3) {$3$};
    \node[left=2pt, font=\footnotesize] at (0, \localMinVal) {$-\frac{175}{27}$};
    \node[right=2pt, font=\footnotesize] at (0, \minVal) {$-6\sqrt{2}-8$};
    \node[left=2pt, font=\footnotesize] at (0, \endVal) {$6\sqrt{2}-8$};

    % 関数名ラベル
    \node[blue!85!black, right=2pt, font=\footnotesize] at (2.2, 3.5) {$y=f(\alpha)$};
  \end{tikzpicture}
  \caption{$f(\alpha)=\alpha^3-\alpha^2-5\alpha$（$-2\sqrt{2}\le\alpha\le2\sqrt{2}$）のグラフ}
  \label{fig:f_alpha_graph}
\end{figure}

\newpage

\end{document}
